\section{Interpretation}
\label{sec:interpretation}

The pilot evaluates a controlled mechanism, not the full reference architecture. The interpretation therefore concerns which feature groups and auxiliary objectives changed behavior on the generated task grammar, how those changes interacted with structural shift and missing context, and which failure modes remained visible.

\subsection{Runtime context}

The B condition isolates the contribution of providing ontology, purpose, and world-state features without axiological auxiliary training. Its comparison with A1 measures the value of explicit runtime context under the same proxy capacity. This is a context-use test: it does not identify whether a future model would ground the symbols in external entities or whether the supplied state is correct.

\subsection{Axiological representation}

C1 adds value and relation objectives to the runtime semantic channels. Improvements over B in the generated primary table are consistent with the mechanism hypothesis that explicit value-bearing relations can provide useful supervision for salience and structural action selection. The result is compatible with simpler explanations, including feature separability and regularities in the generator. Relation-permutation, human-validated data, and matched Transformer experiments remain necessary to identify the cause.

\subsection{Normative training}

C2 adds conflict and authority objectives to the C1 condition. Its conflict results test whether a model can preserve a distinction between no conflict, resolvable conflict, and unresolved conflict rather than collapse those labels into action utility. The interpretation is about classification behavior on the synthetic labels. It does not establish a general theory for resolving incompatible norms or selecting legitimate authority.

\subsection{Structural generalization}

The held-out topology split tests a structural transformation rather than a new surface vocabulary alone. The generated structural-OOD results provide evidence about transfer within the scenario grammar. A gain on this split is informative only alongside the topology-overlap audit, relation-permutation controls, capability, calibration, and per-family error analysis. Structural transfer in this pilot is not evidence that a model learned the intended semantics in an open domain.

\subsection{Candidate-action conformance}

The fixed verifier separates candidate formation from enforcement. A model can improve UCR by producing more useful, well-formed candidates, while a verifier can reject a candidate because it is incomplete or because the fixed schema lacks a required distinction. Coverage results show why this separation matters: removing semantic features can reduce action correctness without necessarily increasing rejection. The conformance measure is consequently a deployment-path diagnostic, not a scalar alignment objective.

\subsection{Semantic occlusion and missing context}

Semantic coverage is a useful derived research gap. Let
\begin{equation}
  K = \operatorname{Coverage}(O_N,O_D,W_t,X,P),
  \label{eq:semantic-coverage}
\end{equation}
where $K$ measures whether the normative, domain, world-state, input, and purpose records expose the entities, relations, exceptions, and effects needed for the current decision. The present pilot does not validate this function as a solution; it motivates measuring it. The semantic-occlusion attack demonstrates the corresponding failure mode: a structurally conditioned model may still select a confident action when a relevant patient or relation is absent. Completeness checks, abstention, provenance, and human review are therefore part of the next experiment rather than optional deployment polish.

The broader architectural tradeoff follows from these observations. Explicit structure may improve separability and auditability, but ontology construction, purpose binding, authority resolution, and state completeness become active sources of error. The value of VGA depends on whether those errors are measurable and containable at larger scale.
